\section{Experimental Details and Evaluation Protocol}
\label{sec:experimental-details}

This section gives the complete protocol behind the reported results.  We distinguish
three quantities that were conflated in the shorter submission: the ambient frame
used to generate each Navier--Stokes example, the relative transform sampled during
training, and the relative transform used for the symmetry-shifted test evaluation.
All ``radii'' below are componentwise bounds.  Thus
\(\delta\sim\mathcal U([-r,r]^2)\), not a uniform draw from a Euclidean ball, and
the step used in the LOCO denominator is
\(\epsilon=\max(\lVert\delta\rVert_2,10^{-6})\leq\sqrt{2}r\).  The clamp is
a numerical guard for the tiny region \(\lVert\delta\rVert_2<10^{-6}\), not only
the exact identity.

\subsection{Headline N64 Galilean experiment}

\paragraph{Equation, domain, and numerical solver.}
We solve the unforced vorticity equation
\begin{equation}
  \partial_t\omega + v\mathbin{\cdot}\nabla\omega
  = \nu\Delta\omega,
  \qquad
  v=(\partial_y\psi,-\partial_x\psi)+c,
  \qquad
  \Delta\psi=\omega,
\end{equation}
on the periodic square \([0,1)^2\).  This sign convention is the one used by the
released solver.  The solution is represented on a \(64\times64\) uniform grid and
advanced from \(t=0\) to \(T=0.5\) with a pseudo-spectral classical fourth-order
Runge--Kutta method, time step \(10^{-3}\) (500 steps), viscosity
\(\nu=10^{-3}\), tensor-product \(2/3\)-rule de-aliasing of the nonlinear term,
and no forcing.  Solver generation uses batches of eight only to limit memory; it
does not change the numerical method.

\paragraph{Initial conditions and ambient frames.}
For every example, we draw independent standard Gaussian noise \(z\) on the grid,
apply the Fourier filter
\(\widehat\omega_0(k)=\widehat z(k)\exp(-\lVert k\rVert_2^2/8)\), subtract the
spatial mean, divide by the sample spatial standard deviation plus \(10^{-6}\),
and use amplitude one.  We then draw the constant ambient-velocity components
independently as \(c_x,c_y\sim\mathcal U[-0.5,0.5]\).  The input channels are
\((\omega_0,c_x,c_y)\in\mathbb R^{64\times64\times3}\), with the two scalar
velocities broadcast over the grid; the target is the terminal vorticity
\(\omega(T)\in\mathbb R^{64\times64\times1}\).  All stored tensors are
\texttt{float32}.  We apply no model-side input or target normalization, discard no
examples, and use no additional data cleaning or clipping.

\begin{table}[t]
\centering
\small
\caption{Complete data-generation parameters for the headline boosted N64
Navier--Stokes dataset.  Generation seed offsets separate the three independently
generated splits.}
\label{tab:n64-dataset-details}
\begin{tabular}{@{}lp{10.5cm}@{}}
\toprule
Item & Value \\
\midrule
Domain / grid & periodic \([0,1)^2\), \(64\times64\) \\
Train / validation / test & 1,024 / 128 / 256 \\
Dataset seed & 31; split offsets 0 / 100,000 / 200,000 \\
Final time / time step & \(T=0.5\) / \(\Delta t=10^{-3}\) \\
Viscosity / forcing & \(10^{-3}\) / none \\
Integrator & pseudo-spectral RK4, 500 steps \\
Aliasing control & tensor-product \(2/3\) filter \\
Initial-condition filter & \(\exp(-\lVert k\rVert_2^2/8)\) \\
IC centering / scaling & zero spatial mean; SD \(+10^{-6}\); amplitude 1 \\
Ambient velocity & \(c_x,c_y\stackrel{\mathrm{iid}}{\sim}\mathcal U[-0.5,0.5]\) \\
Generator batch / precision & 8 / \texttt{float32} \\
Input / target shapes & \(64\times64\times3\) / \(64\times64\times1\) \\
Processed SHA-256 & \texttt{defc9703dbc41494fabf9}\allowbreak\texttt{fcfe9e3408eb2e1c38908}\allowbreak\texttt{14e47fe03d0f6b849e3f50} \\
Config fingerprint & \texttt{b2b5dec2dd3049ff49107}\allowbreak\texttt{1674c34042c89f8c1a855}\allowbreak\texttt{f86cbb199614bfafed0cdd} \\
\bottomrule
\end{tabular}
\end{table}

\paragraph{Splits and labeled subsets.}
The train, validation, and test fields are generated independently and stored before
model fitting.  Validation and test splits are always used in full.  For a training
fraction \(q<1\), a run with seed \(s\) retains the first
\(\max(1,\operatorname{round}(1024q))\) indices of a \texttt{randperm} generated
with seed \(s\).  The resulting 1\%, 2\%, 5\%, and 10\% subsets therefore contain
10, 20, 51, and 102 labeled examples.  Methods with the same seed and fraction use
the same labeled subset.  Consequently, variation across training seeds includes
both parameter initialization and labeled-subset variation, whereas the generated
dataset itself remains fixed.

\paragraph{FNO2d architecture.}
Every headline row uses the same FNO2d with three input channels and one output
channel, width 48, 12 retained modes on each axis, and four Fourier blocks.  Each
block sums a learned spectral convolution and a \(1\times1\) pointwise convolution
and applies GELU.  A pointwise \(48\to96\to1\) projection with GELU produces the
output.  We use no coordinate channels, spatial padding, normalization layers, or
dropout.  The model has 2,668,609 trainable parameters.  LOCO adds no parameters or
inference-time operation.

\paragraph{DeepONet architecture and backbone-control protocol.}
The Reviewer~2 backbone control uses a fixed-grid DeepONet with the same
\(64\times64\), three-channel inputs and one-channel outputs.  Its branch encoder
has four circularly padded \(3\times3\), stride-2 convolutions with channel widths
\(3\to32\to64\to128\to256\), GELU after every convolution, and a
\(4096\to480\to256\) GELU projection.  We deliberately use no normalization:
normalizing spatially constant channels could erase the two ambient-velocity
inputs.  The trunk receives 48 fixed periodic coordinate features,
\(\{\sin(2\pi kx),\cos(2\pi kx),\sin(2\pi ky),\cos(2\pi ky)\}_{k=1}^{12}\),
and applies a \(48\to256\to256\to256\to256\) MLP with GELU after the first
three linear maps.  A rank-256 branch--trunk inner product, divided by
\(\sqrt{256}\), plus one learned scalar bias produces each output-grid value.
The model has 2,688,033 trainable parameters, 0.73\% more than the FNO2d.  The
circular padding and periodic query features encode the domain topology; there
is no spectral convolution and no hard-coded Galilean input--output action.  The
sensor and query grids are fixed at N64, so this experiment does not test
resolution transfer.

We reproduce all four headline training methods with seeds
\(\{23,31,47,59,71\}\) and predeclare DeepONet augmentation versus DeepONet
augmentation plus normalized LOCO as the primary paired comparison.  It shares
the dataset, seed-indexed 20-example labeled subsets, transform distribution,
optimization hyperparameters, and validation rule with the FNO protocol.  All four
DeepONet methods use the same DeepONet inference graph and parameter count.  Methods
share the initialization seed and labeled subset, but their full stochastic training
trajectories are not common-random-number coupled: the schedules restart the
20-example loader a different number of times, and the LOCO method consumes an
additional independent transform draw per update.  Validation and test transform
draws are seed-matched across methods.  We fixed the
architecture and copied the FNO coefficients \(0.05\) (LOCO only) and \(0.10\)
(augmentation plus LOCO) before completing the five-seed test outcomes; there was
no DeepONet-specific test-set or lambda sweep.

\paragraph{Optimization and checkpointing.}
Training uses \texttt{float32} without automatic mixed precision.  We optimize with
AdamW for 150 epochs using learning rate \(10^{-3}\), weight decay \(10^{-4}\),
\((\beta_1,\beta_2)=(0.9,0.999)\), optimizer \(\epsilon=10^{-8}\), and no AMSGrad.
A cosine-annealing schedule spans all 150 epochs, has zero minimum learning rate,
and uses no warm-up.  We clip the global gradient norm at 1.0.  The shuffled
training loader uses batch size 16, no workers, and retains the final partial batch;
validation and test loaders are not shuffled.  Validation runs every 10 epochs on
all 128 examples.  We do not early-stop and select the checkpoint with the lowest
validation ID mean per-example relative \(L^2\).

The supervised and augmented terms are the implemented unsquared per-example
relative \(L^2\) losses in Section~\ref{sec:method}, with
\(\lambda_{\mathrm{aug}}=1\).  The finite-orbit numerator is a mean square over
spatial/output entries and is divided by \(\epsilon^2+\eta\), with
\(\eta=10^{-6}\).  On every relevant update, augmentation and LOCO draw independent
transforms; within each draw, transform parameters are sampled independently for
every example.  Gradients propagate through both LOCO prediction branches.

\begin{table}[t]
\centering
\small
\caption{Headline N64 training variants.  At 2\% labels, batches cycle through
sizes 16 and 4.  The transformed-example totals below use the actual partial
batches rather than treating every update as a full batch.}
\label{tab:n64-training-details}
\resizebox{\linewidth}{!}{%
\begin{tabular}{@{}lcccccc@{}}
\toprule
Method & Steps/epoch & Updates & Batch forwards & \(\lambda_{\rm orb}\) & Aug. examples & Orbit examples \\
\midrule
FNO & 4 & 600 & 600 & -- & 0 & 0 \\
FNO+LOCO & 4 & 600 & 1,200 & 0.05 & 0 & 6,000 \\
FNO+augmentation & 6 & 900 & 1,800 & -- & 9,000 & 0 \\
FNO+augmentation+LOCO & 4 & 600 & 1,800 & 0.10 & 6,000 & 6,000 \\
\addlinespace
DeepONet & 4 & 600 & 600 & -- & 0 & 0 \\
DeepONet+LOCO & 4 & 600 & 1,200 & 0.05 & 0 & 6,000 \\
DeepONet+augmentation & 6 & 900 & 1,800 & -- & 9,000 & 0 \\
DeepONet+augmentation+LOCO & 4 & 600 & 1,800 & 0.10 & 6,000 & 6,000 \\
\bottomrule
\end{tabular}%
}
\end{table}

The primary fixed-weight comparison uses \(\lambda_{\mathrm{orb}}=0.10\); the
orbit-only row uses 0.05.  The robustness sweep evaluates
\(\lambda_{\mathrm{orb}}\in\{0.01,0.05,0.10,0.20,0.30\}\), and we identify 0.30
explicitly when reporting the post-hoc tuned row.  Unless a table says otherwise,
five-seed results use seeds \(\{23,31,47,59,71\}\), while compact controls use
\(\{23,31,47\}\).

\paragraph{Galilean training and transformed-test distributions.}
For a relative boost \(\delta=(\delta_x,\delta_y)\), the input action adds
\(\delta\) to the two ambient-velocity channels and the output action applies the
exact periodic spectral shift \(\omega_T(x)\mapsto\omega_T(x-\delta T)\).
During ordinary training and test evaluation,
\(\delta_x,\delta_y\stackrel{\mathrm{iid}}{\sim}\mathcal U[-0.25,0.25]\).
The severity study instead uses component bounds
\(r\in\{0.10,0.20,0.25,0.35,0.50\}\).  These are additional relative boosts of
the same held-out examples, not a separately simulated test split; because the
original ambient frame is in \([-0.5,0.5]^2\), the transformed observable frame at
\(r=0.50\) can extend to \([-1,1]^2\), although not every draw lies outside the
original componentwise support.  We therefore use ``symmetry-shifted OOD'' in this
specific sense.

\paragraph{Evaluation, uncertainty, and latency.}
For each checkpoint, ID relative \(L^2\) averages one per-example ratio over all
256 test examples.  Orbit OOD error and relative equivariance defect average over
four independent transforms per example (1,024 transformed pairs); the standard
test transform RNG seed is \(s+200{,}000\).  The severity sweep also uses four
draws per example, with seed \(s+700{,}000\).  Tables report the arithmetic mean and
sample standard deviation (\texttt{ddof}=1) across seeds.  Paired comparisons use
matched-seed differences and two-sided 95\% Student-\(t\) intervals with
\(n-1\) degrees of freedom.

Latency is measured on the first 16-example test batch using five unmeasured warm-up
forwards followed by 20 timed forwards; milliseconds per sample is elapsed time
divided by \(20\times16\).  The retained headline manifests record CPU execution,
Windows platform string \texttt{Windows-10-10.0.26100-SP0}, Python 3.11.6, and
PyTorch 2.12.0+cpu.  The host currently identifies its processor as an 11th Gen
Intel Core i7-1185G7 at 3.00 GHz; the original manifests did not record the CPU model,
RAM, BLAS thread count, or peak memory, so latency should be interpreted as
host-specific.  Deterministic kernels were not forced
(\texttt{deterministic=false}), but every run stores its resolved config, config
hash, Git commit, dataset SHA-256, checkpoint, metrics, and initial/final Torch RNG
states.

The DeepONet campaign uses the same host and software versions and records
\texttt{OMP\_NUM\_THREADS=1} and \texttt{MKL\_NUM\_THREADS=1}.  Its seed shards
were executed concurrently as uninterrupted direct runs; checkpoint-chunk resume
was not used.  The recorded run wall times and latency samples
are audit fields rather than controlled cross-method or cross-backbone efficiency
comparisons.  The unchanged-inference statement follows directly from using the
identical DeepONet graph and parameter count with and without LOCO.

\subsection{Secondary datasets and model hyperparameters}

Table~\ref{tab:secondary-dataset-details} records the complete preparation of the
secondary benchmarks.  These experiments are reported as scope and mechanism
checks; the headline claim remains the boosted N64 setting above.

\begin{table*}[t]
\centering
\scriptsize
\caption{Secondary dataset generation, curation, and preprocessing.  All processed
tensors are \texttt{float32}.}
\label{tab:secondary-dataset-details}
\resizebox{\textwidth}{!}{%
\begin{tabular}{@{}p{2.1cm}p{3.1cm}p{2.5cm}p{5.5cm}p{4.6cm}@{}}
\toprule
Dataset & Domain, resolution, splits & Evolution parameters & Initial data or source & Preparation and processed SHA-256 \\
\midrule
1D viscous Burgers & Periodic \([0,1)\), \(N=128\); 2,048/256/512 & \(T=0.5\), \(\Delta t=10^{-3}\), \(\nu=0.01\); pseudo-spectral RK4, 500 steps, \(2/3\) de-aliasing; generator batch 64 & Modes \(k=1,\ldots,8\), independent Gaussian sine/cosine coefficients scaled by \(k^{-1.5}\); zero mean, unit SD, amplitude 0.8, mean 0; dataset seed 23 and split offsets 0/100,000/200,000 & No further normalization or exclusions; shapes \(128\times1\to128\times1\); \texttt{1753fa2e5ca4fb474f93fc8cdef5732d}\allowbreak\texttt{aea94b7d36d7537113b38f9e3d029758} \\
\addlinespace
Unboosted N64 Navier--Stokes (D4) & Periodic \([0,1)^2\), \(64^2\); 1,024/128/256 & Same \(T\), \(\Delta t\), \(\nu\), solver, de-aliasing, and generator batch as the headline data; no ambient velocity & Same filtered Gaussian vorticity law (smoothness 8, amplitude 1); dataset seed 31 and split offsets 0/100,000/200,000 & One input and one target vorticity channel; no further normalization; \texttt{defade883926ad6ba6ef469bf1e6ba69}\allowbreak\texttt{b3f7d57e743e4f53c83a05acb88a8f7e} \\
\addlinespace
1D Dirichlet heat diagnostic & \([0,1]\), \(N=64\) including endpoints; 512/128/128 & \(u_t=0.01u_{xx}\), \(T=0.1\); exact eight-mode sine-series decay \(e^{-0.01(\pi k)^2T}\) & Eight sine modes with independent Gaussian coefficients scaled by \(1/k\), samplewise unit-SD normalization, amplitude 1; split seeds 23/34/45 & Homogeneous Dirichlet endpoints; no further normalization; shapes \(64\times1\to64\times1\); \texttt{7427ec21304ae8879e700653e65ee239}\allowbreak\texttt{3416b2daa8bc4b27756c1cae1ff09b2d} \\
\addlinespace
rMD17 ethanol \citep{christensen2020rmd17} & Fixed nine-atom molecule; 500/1,000/1,000 and 1,000/1,000/1,000 & Static force regression; no PDE solve & Released \texttt{rmd17\_ethanol.npz}; raw SHA-256 \texttt{0824d78c687b8ede25f7162c54c5c643}\allowbreak\texttt{4cfe37f9f859ee8f223c60379d7d5cf9}; historical custom selection with metadata seed 2027, rather than the official split CSVs & Subtract each conformation's arithmetic coordinate centroid; append nuclear charge \(Z/10\); targets are unnormalized forces. Exact cached 500-label tensor is preserved in the artifact with raw \texttt{old\_indices}; SHA-256 \texttt{fb5515b3355651459c2ad9078e7cfb5}\allowbreak\texttt{bb63dce96e55e7c36956cbe9ff25e5472}. Coordinates in \AA, forces in kcal mol\(^{-1}\)\AA\(^{-1}\) \\
\bottomrule
\end{tabular}%
}
\end{table*}

\begin{table*}[t]
\centering
\scriptsize
\caption{Model and optimization hyperparameters for every reported benchmark.
Settings not repeated in the last column use the AdamW/cosine defaults stated in
the main protocol.}
\label{tab:all-model-training-details}
\resizebox{\textwidth}{!}{%
\begin{tabular}{@{}p{2.3cm}p{6.4cm}p{2.0cm}p{2.6cm}p{5.0cm}@{}}
\toprule
Benchmark & Architecture & Parameters & Epochs / batch & Benchmark-specific training and transforms \\
\midrule
Boosted N64 & FNO2d, \(3\to1\), width 48, modes \(12\times12\), depth 4, no grid, GELU & 2,668,609 & 150 / 16 & 2\% headline has 20 labels and 4 steps/epoch; augmentation compute match uses 6; \(\lambda_{\rm aug}=1\), fixed \(\lambda_{\rm orb}=0.10\), \(\eta=10^{-6}\) \\
\addlinespace
Boosted N64 / DeepONet & Circular CNN branch \(3\to32\to64\to128\to256\), \(4096\to480\to256\) head; 12-harmonic periodic coordinate trunk, width/rank 256 & 2,688,033 & 150 / 16 & Same 20-label subsets, seeds, optimizer, transform law, and 4/4/6/4 step schedules as the FNO four-method protocol; fixed \(\lambda_{\rm orb}=0.10\) for the primary Aug.+LOCO row \\
\addlinespace
Burgers & FNO1d, \(1\to1\), width 64, 16 modes, depth 4, no grid, GELU & 287,361 & 150 / 32 & Full 2,048 labels, 64 steps/epoch; seeds 23/31/47; orbit-only \(\lambda=0.05\), Aug.+LOCO \(\lambda=0.25\) \\
\addlinespace
Dirichlet heat & FNO1d, \(1\to1\), width 32, 12 modes, depth 4, coordinate grid, GELU & 55,649 & 80 / 32 & 10\%=51 labels, 2 steps/epoch; seeds 23/31/47; \(\lambda_{\rm aug}=1\), \(\lambda_{\rm orb}=0.05\), \(\eta=10^{-6}\) \\
\addlinespace
rMD17 ethanol & Flatten \(9\times4\) input; four width-128 SiLU linear layers; 27-output linear head reshaped to \(9\times3\) & 57,755 & 150 / 32 & Reported 500-label comparison: Aug. 24 and Aug.+LOCO 16 steps/epoch (48 batch forwards/epoch), seeds 23/31/47/59/71; validation every 25 epochs; \(\lambda_{\rm aug}=1\), \(\lambda_{\rm orb}=0.03\), \(\eta=10^{-3}\); step normalization enabled \\
\addlinespace
D4 FNO & FNO2d, \(1\to1\), width 48, modes \(12\times12\), depth 4, no grid, GELU & 2,668,513 & 150 / 16 & 1\%=10 labels, 4 steps/epoch; seeds 23/31/47/59/71; Aug.+LOCO \(\lambda=0.20\), \(\eta=10^{-6}\) \\
\addlinespace
D4 G-FNO & D4 group FNO, pseudoscalar input/output, width 8, modes \(6\times6\), depth 4, no grid & 350,280 & 150 / 16 & Supervised only; same 10-label subsets, 4 steps/epoch, and five seeds; architectural comparator is not parameter- or mode-matched \\
\bottomrule
\end{tabular}%
}
\end{table*}

\paragraph{Burgers transform mixture.}
The reported Burgers orbit metrics are not Galilean-only.  For each minibatch, we
choose either a periodic translation (probability 0.5, shift
\(s\sim\mathcal U[-0.25,0.25]\)) or a Galilean boost (probability 0.5,
\(b\sim\mathcal U[-0.35,0.35]\)); parameters within the selected family are then
independent across examples.  The Galilean action is
\(u_0(x)\mapsto u_0(x)+b\) and
\(u_T(x)\mapsto u_T(x-bT)+b\).  Severity multipliers
\(1,1.5,2,3\) produce translation bounds
\(0.25,0.375,0.50,0.75\) and boost bounds
\(0.35,0.525,0.70,1.05\).  Standard evaluations use four draws per example;
the severity analysis uses eight.

\paragraph{Molecular rigid motions.}
For rMD17, each transform draws an axis by normalizing a three-dimensional
standard Gaussian vector, an angle uniformly in \([-\pi/2,\pi/2]\), and each
translation component uniformly in \([-1,1]\)~\AA.  Coordinates rotate and
translate, whereas force vectors rotate only.  For this product action we disclose
the implemented composite step
\(\epsilon=\max(|\theta|+\lVert t\rVert_2,10^{-6})\), with angle in radians and
translation in \AA.  This is a disclosed composite step weighting that mixes the
configured angular and translational units; it is not a canonical, unit-free
\(\mathrm{SE}(3)\) norm.  The reported LOCO row divides by
\(\epsilon^2+10^{-3}\) and uses
\(\lambda_{\rm orb}=0.03\), a rounded approximate average-scale match to the
earlier raw coefficient 0.01 under this transform law.  We use 10 latency warm-ups
and 50 repeats.
For the reported 500-label comparison, augmentation performs 24 updates per epoch
with two model evaluations per update, whereas Aug.+LOCO performs 16 updates with
three evaluations per update.  Both therefore use 7,200 batch-level model
evaluations over 150 epochs, although their optimizer and backward-pass counts
differ.
The historical cache predates the current rMD17 generator implementation.  Although
its metadata records split seed 2027, a fresh \texttt{torch.randperm} does not
reconstruct the same selection.  The exact 500/1,000/1,000 tensor used here is
therefore included as the authoritative evidence input, and its retained raw
\texttt{old\_indices} make the selected conformations auditable against the hashed
NPZ.  The released verification script matches all 2,500 cached conformations to
the raw NPZ with zero \texttt{float32} input and force discrepancy.  The older
1,000-label cache has different validation/test slices and is not
used as a controlled label-scaling comparison.

\paragraph{Finite D4 action.}
For the D4 experiment, a rotation index is uniform on
\(\{0,1,2,3\}\) and reflection is Bernoulli\((0.5)\), giving all eight elements
uniformly.  Vorticity is treated as a pseudoscalar and changes sign under reflection.
We set \(\epsilon=1\), so this is finite-group consistency rather than a local
Lie-algebra step.  Augmentation and Aug.+LOCO both use four updates per epoch; their
forward-pass compute is therefore not matched.

\paragraph{Boundary diagnostic.}
The heat-equation shift is sampled as \(s\sim\mathcal U[-0.1,0.1]\) and implemented
by zero-padded linear interpolation.  The masked variant evaluates only source
coordinates in \([0.05,0.95]\).  Because a fixed Dirichlet problem is not globally
translation-equivariant and boundary influence propagates into the interior, this
experiment is an approximate boundary-mask diagnostic, not an exact-symmetry OOD
benchmark.  Its transformed target is the shifted cached output, not a newly solved
heat equation for the shifted input.  To avoid comparing checkpoints on different
scoring supports, we re-evaluate all three training variants with the same transform
draws under both a full-domain score and a common valid-overlap score with margin
0.05.  Masking restricts where this approximate constraint is scored; it does not
restore translation equivariance of the fixed-boundary heat equation.

\subsection{Additional N64 controls}

\paragraph{Step-size and tangent controls.}
At 2\% labels, the denominator sensitivity uses
\(\eta\in\{10^{-8},10^{-6},10^{-4}\}\) at component radius 0.25, while the
step-size sensitivity uses radii \(\{0.125,0.25,0.50\}\) at \(\eta=10^{-6}\);
each uses seeds 23, 31, and 47.  The tangent-propagation comparison uses the same 20
labels, four steps per epoch, all five headline seeds, and
\(\lambda_{\mathrm{tangent}}=0.10\).  Its input tangent changes only the two ambient
velocity channels, and its output tangent is
\(-T(\delta_x\partial_x+\delta_y\partial_y)\mathcal G(a)\).  A differentiable JVP
is computed with graph construction enabled, and the tangent mean square is divided
by \(\lVert\delta\rVert_2^2+10^{-6}\).  Augmentation+tangent draws its supervised
augmentation independently.

The plain finite-transform control uses the same data, architecture, transform
range, augmentation, four-update schedule, and five headline seeds
\(\{23,31,47,59,71\}\) as normalized LOCO, but minimizes the raw finite
equivariance MSE.  Since
\(\delta_x,\delta_y\stackrel{\mathrm{iid}}{\sim}\mathcal U[-0.25,0.25]\),
\(\mathbb E\lVert\delta\rVert_2^2=2(0.25)^2/3=1/24\).  We therefore set its
raw-loss weight to 2.4, an approximate leading-order average-scale match to the
normalized \(\lambda_{\mathrm{orb}}=0.10\) objective.  The square sampling law
couples direction and norm, so this is not an exact equivalence, nor an assertion
that the two nonlinear losses have identical realized magnitudes.

\paragraph{Semi-supervised joint training.}
The joint 2\%-label experiment uses supervised and augmented batches from the
20-example labeled subset and separate LOCO batches from the full 1,024-example
training pool.  It performs four combined updates per epoch, uses full-pool orbit
batch size 16, and therefore samples 64 (not all 1,024) full-pool orbit examples per
epoch; \(\lambda_{\mathrm{orb}}=0.05\) and seeds are 23, 31, and 47.

\paragraph{Unlabeled target adaptation.}
The target-adaptation table is transductive: the same 256 test inputs are used without
labels during adaptation and with labels for before/after evaluation.  Starting from
the 2\%-label augmentation or fixed-weight Aug.+LOCO checkpoint, we update only the
FNO projection layers (4,801 parameters) for five epochs and 80 updates.  We use
AdamW with learning rate \(10^{-5}\), zero weight decay, no scheduler, batch size 16,
and gradient clip 0.1.  The objective combines orbit loss with a unit-weight relative
prediction-preservation term; parameter anchoring has weight zero.  Target and orbit
boosts are independently redrawn at component radius 0.35 or 0.50.  Evaluation uses
four transforms per example and seeds 23, 31, and 47.

\paragraph{Canonicalization diagnostics.}
The metric called oracle paired-source error compares
\(T_g\mathcal G(a)\) with \(T_gu\); it does not evaluate
\(\mathcal G(T_ga)\) and is therefore a diagnostic rather than a new model.
Observable canonicalization is also evaluation-only: it reads the total ambient
velocity from the input, sets those channels to zero, evaluates the FNO once, and
shifts the result by \(cT\).  It changes preprocessing but neither retrains the
checkpoint nor changes the number of model calls.
